% Generated draft: requires mathematical and source review.
\documentclass[11pt]{article}
\usepackage[margin=25mm]{geometry}
\usepackage{fontspec}
\setmainfont{Latin Modern Roman}
\usepackage{amsmath,amssymb,mathtools,unicode-math}
\setmathfont{Latin Modern Math}
\usepackage{enumitem,needspace,xurl,hyperref}
\setlength{\parindent}{0pt}
\setlength{\parskip}{5pt}
\setlength{\emergencystretch}{3em}
\newcommand{\N}{\mathbb{N}}
\newcommand{\Z}{\mathbb{Z}}
\newcommand{\Q}{\mathbb{Q}}
\newcommand{\R}{\mathbb{R}}
\newcommand{\C}{\mathbb{C}}
\newcommand{\F}{\mathbb{F}}
\newcommand{\A}{\mathbb{A}}
\newcommand{\PP}{\mathbb{P}}
\newcommand{\E}{\mathbb{E}}
\newcommand{\eps}{\varepsilon}
\newcommand{\dd}{\,\mathrm{d}}
\begin{document}
\section*{UnsolvedMath ID: 30002891}
\textbf{Source problem code:} catalogue entry
\subsection{Definitions and mathematical statement}
For a finite-dimensional real or complex matrix \(A\), its trace norm is
\[
\|A\|_1=\operatorname{tr}\bigl((A^*A)^{1/2}\bigr),
\]
where \(A^*\) denotes the transpose or conjugate transpose, respectively. The supplied record does not specify the matrix dimensions, scalar field, admissible parameter class, observation scheme, or sampling distribution; these must be supplied before a particular statistical problem can be stated.

To explain the confidence-region terminology conditionally, let \(\Theta\) be a specified matrix parameter space, let \(X\) denote observed data with law \(P_\theta\) for \(\theta\in\Theta\), and fix a confidence error probability \(\alpha\in(0,1)\). A measurable random set \(C_\alpha(X)\subseteq\Theta\) has frequentist coverage at least \(1-\alpha\) over this model if
\[
\forall\theta\in\Theta,\qquad
P_\theta\{\theta\in C_\alpha(X)\}\geq 1-\alpha.
\]
One possible trace-norm construction is
\[
C_\alpha(X)=\{\vartheta\in\Theta:
\|\vartheta-\widehat\theta(X)\|_1\leq r_\alpha(X)\},
\]
where \(\widehat\theta\) is an estimator and \(r_\alpha(X)\geq0\) is a measurable radius. This is a conventional interpretation of a trace-norm confidence region, not a construction or coverage requirement explicitly specified by the source. The record does not determine whether its word ``interval'' means such a region or an interval for a scalar functional.

\textbf{Source research agenda.} Develop practical trace-norm confidence intervals and trace-norm certificates; characterize confidence intervals in other norms; and analyze alternative construction methods, including Bayesian and bootstrap procedures.

The supplied agenda leaves the meaning of ``certificate,'' the other norms, the required accuracy and computational guarantees, and the intended finite-sample or asymptotic coverage unspecified. It also does not specify whether Bayesian constructions are to provide posterior credibility, frequentist coverage, or both. Consequently, the agenda cannot be strengthened into a uniquely determined theorem or quantified existence claim from this record alone.
\subsection{Short English statement}
The problem asks for practical uncertainty quantification in matrix estimation using the trace norm and other norms, including certificates and Bayesian or bootstrap constructions. The statistical model and required guarantees are unspecified, so the displayed coverage condition is a conditional interpretation requiring human review. The supplied status is partially solved; it is not independently verified here.
\subsection{Sources}
\begin{itemize}
\item Supplied local UnsolvedMath registry record, permanent ID \(30002891\), title ``Practical Confidence Regions for Matrix Estimation,'' exact source problem code \texttt{catalogue entry}. No individual proposer is supplied.
\item Supplied report citation: \emph{Probabilistic Techniques in Modern Statistics} (2016), Oberwolfach Reports, DOI 10.4171/owr/2015/26, pp.\ 1449--1498. \url{https://doi.org/10.4171/owr/2015/26}. The record attributes the agenda to workshop open-problem discussions and separately lists proposed year 2015; the date relationship is unresolved.
\item Supplied literature reference: \emph{Projected Least-Squares Quantum Process Tomography}. \url{https://quantum-journal.org/papers/q-2022-10-20-844/}.
\item Supplied literature reference: Faist--Renner, \emph{Physical Review Letters}. \url{https://doi.org/10.1103/PhysRevLett.117.010404}.
\item Supplied literature reference: Oberwolfach report 26. \url{https://doi.org/10.4171/owr/2015/26}.
\item Literature assessment dated exactly \texttt{2026-08-21}: the supplied curation labels the agenda partially solved and reports later practical and nonasymptotic confidence regions in several trace or related norms, while broader certificates, other norms, and bootstrap or Bayesian constructions leave the overall agenda incomplete. This is supplied dated curation, not an independently verified literature assessment; its time-sensitive status and cited claims require independent verification.
\end{itemize}
\end{document}
